\documentclass[carta,firma,final]{lafrox}



\datoslafrox{
  subtitulo    = {Formulario T-9 --- Metodología de gestión de proyecto},
  alcance      = {Formulario T-9},
  formulario   = {Formulario T-9},
  nomenclatura = {LAFROX-Formulario-T-9.pdf},
}

\begin{document}

\portadalafrox
\preliminareslafrox

\renewcommand{\thetable}{T-9.\arabic{table}}
\renewcommand{\thefigure}{T-9.\arabic{figure}}


\chapter{Gestión de proyecto}

Para que el desarrollo del producto final sea un éxito total, 
se debe por necesidad seguir metodologías asociadas a la gestión de proyecto y a su vez, 
metodologías para el desarrollo de software, estas metodolgías seran las siguientes

\section{Metodología de gestión de proyecto}

La gestión del proyecto aplica la Guía del PMBOK, sexta edición (PMI, 2017), adaptada a un proyecto de 56 meses con dos etapas, solapamiento entre ellas y una operación que no puede detenerse; se usan los procesos de integración, alcance, cronograma, calidad, recursos, comunicaciones, riesgos, adquisiciones e interesados, con la profundidad que exige cada uno. El alcance y el cronograma tienen líneas base formales, porque los hitos del Formulario E-25 son fechas fijas y se aceptan por acta. El avance se mide con valor ganado y no con porcentajes declarados (RT-19.07). El flujo diario del equipo se gestiona con un tablero Kanban, que no reemplaza la línea base ni la aceptación formal de los entregables.

La planificación define los entregables, sus responsables, los recursos, los plazos y las dependencias entre actividades, considerando las restricciones de la operación. La línea base de alcance (paquete 1.2.1, H1) y la de cronograma (1.3.3) sirven para medir el avance y evaluar las desviaciones.

El proyecto se ejecuta mediante entregas incrementales dentro de las dos etapas contractuales. La distribución propuesta prioriza primero la calidad y disponibilidad de la información operacional, y luego las capacidades que dependen de esos datos para optimizar y analizar la operación. En la Etapa 1 se implementan las capacidades de preventa, recepción, bodega, preparación y cross-docking; la trazabilidad de los lotes y de la cadena de frío; la planificación de rutas basada en el conocimiento del planificador, con el acuse de recibo de los conductores; y el manejo de efectivo y los retiros de los pedidos. En la Etapa 2 se implementan los portales y el intercambio electrónico con los clientes y, por último, el costo de servir.

La planificación respeta el cronograma contractual: desarrollo de la Etapa 1 entre los meses 1 y 12, marcha blanca entre los meses 13 y 15 y paso a producción en el mes 16. El desarrollo de la Etapa 2 ocupa los meses 13 a 18, la marcha blanca los meses 19 y 20 y el paso a producción el mes 21. El trabajo de la Etapa 2 se coordina con la marcha blanca y la estabilización de la Etapa 1, resguardando la continuidad operacional y la integridad de los datos compartidos.

El tablero Kanban muestra el trabajo pendiente, en ejecución, en revisión y terminado, con responsables y bloqueos, y fija límites de trabajo en curso. Cada tarea tiene una definición de terminado y, cuando corresponde, evidencia asociada a los criterios de aceptación del entregable al que contribuye. La gestión del proyecto se complementa con los procesos de los apartados siguientes.

\subsection{Gestión de interesados}

La gestión de interesados identifica a las personas, grupos y organizaciones que pueden influir en el proyecto o verse afectados por sus resultados. Para cada interesado se registran sus necesidades, expectativas, nivel de influencia, impacto potencial y disposición frente a los cambios. El análisis distingue quiénes participan en decisiones, quiénes aportan conocimiento operacional, quiénes validan entregables y quiénes son afectados o ejercen supervisión externa.

El jefe de proyecto mantiene actualizado el registro de interesados (paquete 1.4.1) y planifica su involucramiento de acuerdo con el alcance de cada entrega. La participación puede incluir levantamiento de necesidades, revisión de procesos, validación de criterios, pruebas y retroalimentación. Las inquietudes, resistencias, acuerdos y compromisos relevantes se registran con su responsable y seguimiento; los conflictos que no puedan resolverse en la instancia de trabajo se elevan a la autoridad correspondiente.

\parte{subdocumento_6/tablas/interesados}

La tabla resume los principales grupos y roles identificados, sus intereses y la contribución esperada. La participación concreta se acuerda según las responsabilidades de cada interesado y las actividades de cada entrega; la inclusión en la tabla no implica que todos participen en todas las decisiones ni que tengan atribuciones de aprobación.

\subsection{Gestión de comunicaciones}

La gestión de las comunicaciones mantiene informados a los responsables de Puelche y a los demás interesados sobre los avances de las entregas que les competen, sostiene una comprensión común del estado del proyecto, facilita la toma de decisiones y comunica oportunamente las situaciones que puedan afectar los compromisos acordados.

Las reuniones convocan a las personas necesarias para analizar y resolver sus asuntos. Los acuerdos se registran en actas que identifican las decisiones, las actividades pendientes, sus responsables y los plazos; las actas de los comités se levantan dentro de los dos días hábiles siguientes (RT-19.09). La documentación vigente, los entregables, las actas y los registros de riesgos y de cambios se mantienen en el espacio colaborativo accesible al CLIENTE (RT-19.05).

Antes de cada entrega incremental se informa a los interesados pertinentes sobre los cambios previstos, las actividades de validación y el apoyo necesario para la puesta en funcionamiento. Las observaciones se registran y reciben respuesta, distinguiendo las que pueden atenderse dentro de lo acordado de las que requieren una solicitud de cambio.

\parte{subdocumento_6/tablas/comunicaciones}

\subsection{Gestión de adquisiciones}

La gestión de adquisiciones planifica, contrata y controla los bienes y servicios que requiere la solución: contratos, órdenes de compra, acuerdos con terceros y acuerdos de nivel de servicio. Cada adquisición tiene un responsable, una fecha de necesidad derivada del cronograma, una dependencia con los paquetes que la usan y una evidencia de recepción.

El CLIENTE compra el equipamiento de terreno y el de la sala técnica conforme a la especificación de cantidades y características que prepara LafroX (Caso 02, capítulo 11; SD3 E-09; paquetes 5.1.1 y 5.1.2). Esa compra no traslada al CLIENTE la instalación: LafroX habilita el recinto técnico en el espacio que proporciona el CLIENTE, instala y configura los equipos (fase 6 de la EDT) y contrata los servicios que le corresponden. Las licencias de terceros se constituyen a nombre del CLIENTE (5.2.3).

\parte{subdocumento_6/tablas/adquisiciones}

Cada fila tiene en el registro de adquisiciones su estado, su proveedor y su fecha comprometida. Un atraso que amenace el H3 o una ola se escala al Comité Ejecutivo con su análisis de impacto.

\subsection{Gestión de integración}

La gestión de integración mantiene alineados los componentes del proyecto para que las decisiones sobre alcance, cronograma, costos, recursos, calidad, riesgos y adquisiciones se analicen de manera coordinada. El jefe de proyecto consolida estos elementos en el plan para la dirección del proyecto, lo comunica a los responsables y lo actualiza cuando se aprueban cambios.

Durante la ejecución, el jefe de proyecto coordina las actividades y sus dependencias, revisa el avance de los entregables y gestiona los impedimentos que afectan a más de un equipo o área. Las decisiones, acuerdos, riesgos transversales y asuntos que requieran escalamiento quedan registrados con su responsable, resolución y fecha de seguimiento.

Toda solicitud de cambio que pueda afectar el alcance, los plazos, los costos o los criterios de aceptación se documenta y evalúa antes de implementarla. El análisis considera sus efectos en el plan, los entregables relacionados, los riesgos y la operación. El jefe de proyecto eleva la solicitud a la autoridad correspondiente de Puelche y comunica la decisión a las personas afectadas. Solo los cambios aprobados se incorporan a la planificación vigente. Las acciones correctivas destinadas a cumplir los compromisos aprobados se registran y supervisan; si modifican esos compromisos, se tramitan además como solicitudes de cambio.

Al finalizar cada entrega se comprueba que sus resultados satisfagan los criterios de aceptación y cuenten con la evidencia requerida. Las observaciones pendientes se registran con responsable y tratamiento acordado, y se documentan las lecciones aprendidas.

\subsection{Control del alcance, del cronograma y del valor ganado}

El alcance se controla contra la línea base y la matriz de trazabilidad del Formulario T-12: un requerimiento sólo se da por cumplido cuando pasa la prueba asignada en el paquete que lo construye. El cronograma se controla contra la red del Formulario T-15, sección 5, que incluye los diez días hábiles de revisión del CLIENTE antes de cada hito (Art. 18.3).

El avance se mide con valor ganado (RT-19.07). El valor planificado de cada mes es la suma de las horas hombre programadas de los paquetes en curso según el Formulario T-15; el valor ganado acredita las horas de un paquete sólo en hitos de avance definidos de antemano (inicio, revisión interna y aceptación) y no por porcentaje declarado; el costo real son las horas registradas. Con ellos se calculan el índice de desempeño del cronograma (SPI) y el del costo (CPI). Un SPI o un CPI bajo 0,90, o una desviación mayor al 10 %, obligan a presentar un plan de recuperación dentro de cinco días hábiles (paquete 1.8.6).

Las reservas se controlan por separado: la capacidad protegida de la Etapa 1 y las necesidades de contingencia del SD8 sólo se usan con autorización registrada, y su consumo se informa en el mismo informe mensual.

\subsection{Mecanismos de decisión y cadencias de gobierno}

Los comités del Art. 71° de las Bases Administrativas se organizan así:
\begin{itemize}
    \item Reunión de seguimiento semanal (paquete 1.3.5), con el jefe de proyecto y líderes de frente: revisa avance, bloqueos y datos para los comités; no reemplaza a ningún comité.
    \item Comité de Proyecto quincenal (paquete 1.8.3), con jefe de proyecto, líderes y Contraparte Técnica: revisa avance de entregables y el registro vivo de riesgos (RT-19.04).
    \item Comité Ejecutivo mensual (paquete 1.8.2), con gerencias del CLIENTE y dirección de LafroX: resuelve escalamientos, cambios de alcance o plazo y uso de reservas.
    \item Comité de Arquitectura mensual (paquete 1.8.4), con Arquitecto de Solución, Seguridad y TI del CLIENTE: decide asuntos de arquitectura y mantiene su registro.
    \item Comité de Operación mensual desde el mes 13 hasta el mes 56 (paquete 1.8.5), con Líder de Operación / SRE, mesa de servicio, Operaciones y TI del CLIENTE: revisa niveles de servicio, incidentes y problemas, capacidad de la mesa, pruebas de continuidad y plan de mejoras.
\end{itemize}

El Comité de Operación empieza con la primera marcha blanca, porque desde entonces la solución atiende operaciones reales. Revisa cada mes la disponibilidad de los servicios críticos, los incidentes críticos y altos y sus causas, la atención de la mesa (respuesta antes de 20 segundos, abandono y resolución al primer contacto), las pruebas de restauración y de continuidad y los cambios programados. Sus actas registran acuerdos, responsables y plazos (RT-19.09).


\end{document}